\begin{table}[!h]
\centering
\caption{\label{tab:supp_h3_equivalence}H3, the registered prediction that the attribution penalty does not depend on the kind of image (figurative or non-figurative): the attribution-by-image-kind interaction, in rating points, tested against the registered bound of plus or minus 0.20 rating points on each dimension and under each of two random-effects structures. The intercepts-only structure gives each participant and each image its own baseline; the maximal structure also lets the image-kind difference vary from participant to participant and the attribution penalty vary from image to image, which widens the interval on the interaction and so makes a prediction of no effect harder, not easier, to confirm. Each row reads two tests together: the equivalence test, two one-sided tests of the estimate against each end of the bound, each at .05, so the interval shown is a 90\% confidence interval and tost\_p is the larger of the two one-sided p-values (equivalent is TRUE when both ends are rejected), and the ordinary test of whether the interaction differs from zero (nhst\_p, nhst\_significant). Not significant and equivalent is supported; significant and not equivalent is rejected; either remaining combination is inconclusive, and the verdict column states which. H3 is claimed only if the verdict is the same under both structures, and a disagreement is reported as inconclusive. A singular fit is one in which the optimiser has driven an estimated variance or correlation to the edge of its range, a sign that the data cannot support that many parameters; it can still carry slope variances above zero, so the fallback below is the registered rule for such a fit and not a finding that the slopes are absent. When the maximal fit is singular the verdict rests on the intercepts-only structure and the table records which structure was used: random\_structure names the structure of each row, fit\_singular flags the singular fit, and h3\_claim states the conclusion for the dimension. The rows at the wider bounds of 0.30 and 0.50 rating points in the same file are descriptive sensitivity, carry no verdict, and are left out here. A passed equivalence test stands on its own, because the two one-sided tests hold the rate of wrongly declaring equivalence at .05 on each side whatever their power; a failed one is read against the simulated power of the equivalence test at the registered bound under a true interaction of zero, deposited in tables/study1\_h3\_equivalence\_power.csv and rendered in the table that follows.}
\centering
\resizebox{\ifdim\width>\linewidth\linewidth\else\width\fi}{!}{
\fontsize{9}{11}\selectfont
\begin{tabular}[t]{llrrrrlllllll}
\toprule
dimension & random\_structure & estimate & se & ci90\_lo & ci90\_hi & tost\_p & equivalent & nhst\_p & nhst\_significant & fit\_singular & verdict & h3\_claim\\
\midrule
creativity & intercepts only: (1 | id) + (1 | stimulus\_id) & -0.081 & 0.095 & -0.237 & 0.075 & 0.105 & FALSE & 0.393 & FALSE & FALSE & inconclusive: cannot separate no effect from a small one & inconclusive: cannot separate no effect from a small one\\
effort & intercepts only: (1 | id) + (1 | stimulus\_id) & -0.020 & 0.096 & -0.178 & 0.137 & 0.030 & TRUE & 0.833 & FALSE & FALSE & supported & supported\\
creativity & maximal: (1 + stimulus\_type | id) + (1 + Type | stimulus\_id) & -0.081 & 0.110 & -0.269 & 0.108 & 0.145 & FALSE & 0.468 & FALSE & TRUE & inconclusive: cannot separate no effect from a small one & inconclusive: cannot separate no effect from a small one\\
effort & maximal: (1 + stimulus\_type | id) + (1 + Type | stimulus\_id) & -0.020 & 0.096 & -0.179 & 0.139 & 0.032 & TRUE & 0.834 & FALSE & TRUE & supported & supported\\
\bottomrule
\end{tabular}}
\end{table}